\subsection{The cyclotomic restriction on a scalar building prime}

\begin{statement}{A.135}{Proved}
{The actual norm of one scalar Kummer ramification prime}
Let \(r\ge1\), put \(m=3^r\), and let \(K\) be a
number field containing \(\mu_m\).
Let \(q\ne3\) split completely in \(K\), and fix a
prime \(\mathfrak q\) above it.
Suppose \(z\in K^\times\) has valuations divisible
by \(m\) at every finite prime other than
\(\mathfrak q\), has valuation prime to three at
\(\mathfrak q\), and satisfies
\[
                        z\in(K_v^\times)^m
                                    \quad\text{for every }v\mid3.
 \tag{A.135.1}
\]
Then
\[
                          q\equiv1\pmod{3^{r+1}}.
 \tag{A.135.2}
\]
In particular, every single-prime scalar building
class in the odd-power construction used for A.130,
with zero local class above three, has this restriction.
It holds even when its initial splitting field does
not contain \(\mu_{3^{r+1}}\).

For a general finite list of distinguished primes
\(\mathfrak q_i\mid q_i\), where each \(q_i\) splits
in \(K\), the same argument gives the actual
necessary condition
\[
              \prod_iq_i^{a_i}\in(\Q_3^\times)^m,
                  \qquad a_i=v_{\mathfrak q_i}(z)\pmod m.
 \tag{A.135.3}
\]
The residues of the exponents determine this class.
This product condition permits cancellation between
different building primes.  For example \(61\) and
\(67\) both split \(\Q(3^{1/3},\mu_3)\), have
\(v_3(q_i-1)=1\), and their product is a cube in
\(\Q_3^\times\).
This numerical comparison asserts the product
condition, not the existence of an additional
global Kummer class with specified prime divisors.
\end{statement}

\begin{proof}
Choose an integer representative \(a\) of
\(v_{\mathfrak q}(z)\bmod m\), so \(3\nmid a\).
For every rational prime \(p\), the actual
valuation of the field norm is
\[
            v_p(N_{K/\Q}z)
                      =\sum_{v\mid p}f(v/p)v(z).
 \tag{A.135.4}
\]
Since \(q\) splits completely, the residue degree
of \(\mathfrak q\) is one.  The hypotheses give
valuation \(a\bmod m\) at \(q\) and zero modulo
\(m\) at every other rational prime.
A rational number with those valuations has the
literal form
\[
                       N_{K/\Q}z=q^a w^m,
                                 \qquad w\in\Q^\times.
 \tag{A.135.5}
\]
The possible negative sign is absorbed in \(w\)
because \(m\) is odd.

At three the norm is the product of the norms
from every completion above three.
Writing \(z=z_v^m\) in each \(K_v\) shows
\[
       N_{K/\Q}z
             =\left(\prod_{v\mid3}N_{K_v/\Q_3}z_v\right)^m
                                  \quad\text{in }\Q_3^\times.
 \tag{A.135.6}
\]
Equations (A.135.5)--(A.135.6) make \(q^a\) an
\(m\)th power in \(\Q_3^\times\).
Choose integers \(A,B\) with \(Aa+Bm=1\).
Then \(q=(q^a)^A(q^B)^m\) is itself an \(m\)th
power.  Thus no unproved implication about the
norm group is needed.

The decomposition
\(\Z_3^\times=\{1,-1\}\times(1+3\Z_3)\)
and the three-adic logarithm identify the image
of the \(3^r\)-power map with
\(\{1,-1\}(1+3^{r+1}\Z_3)\).
Equivalently this follows by induction from
\((1+3^j u)^3\equiv1+3^{j+1}u\bmod3^{j+2}\).
Hence \(q\equiv1\) or \(-1\bmod3^{r+1}\).
Since \(K\) contains \(\mu_3\) and \(q\) splits
in \(K\), its action on \(\mu_3\) is trivial;
thus \(q\equiv1\bmod3\).
Only the positive alternative is possible,
proving (A.135.2).

In the scalar construction the cohomology class
is a Kummer class in \(K^\times/(K^\times)^m\).
Being unramified at a prime away from three
means that its valuation is zero modulo \(m\).
Its single distinguished inertia coordinate is
one modulo \(m\), and its local class at every
place above three is zero.  Thus any representative
satisfies precisely the hypotheses just used.
The conclusion concerns those actual building
classes and does not impose the single-prime
condition on their final product.

For several distinguished primes the same
valuation formula gives
\(N_{K/\Q}z=(\prod_iq_i^{a_i})w^m\);
(A.135.6) proves (A.135.3).
For the explicit comparison, both \(61\) and
\(67\) are one modulo three, and the congruence
\(X^3=3\) has three roots at each prime.
They therefore split the displayed cubic
splitting field.  Their residues modulo nine
are seven and four, and
\(61\cdot67\equiv1\bmod9\).
The unit power calculation above makes their
product a cube in \(\Q_3^\times\), while neither
individual prime is a cube there.
This proves the exact scope of the comparison.
\end{proof}

The retained certificate
\path{code/scalar_cyclotomic_precision_certificate.py}
checks the actual finite unit power maps through
\(r=6\), the split-prime example, the norm-exponent
B\'ezout identity and the multiple-prime cancellation.
Its output is saved alongside it.
